% SLDE-AFT manuscript -- article-class skeleton.
% TODO: swap this preamble for the Springer journal's official LaTeX
% template (e.g. svjour3/sn-jnl) once the target journal is confirmed;
% article class is a placeholder for drafting only.
\documentclass[11pt]{article}
\usepackage[utf8]{inputenc}
\usepackage[T1]{fontenc}
\usepackage{amsmath,amssymb,amsthm}
\usepackage{graphicx}
\usepackage[numbers,sort&compress]{natbib}
\usepackage{booktabs}
\usepackage{hyperref}

\newtheorem{proposition}{Proposition}
\newtheorem{corollary}{Corollary}

\title{SLDE-AFT: A Self-Learning Dual-Source Knowledge Extraction and Automated Fine-Tuning System}
\author{Shashank Kumar Sadashiva \and Ahmed Ebada}
\date{}

\begin{document}
\maketitle

% ------------------------------------------------------------------
\section{Introduction}
% TODO: not yet written. Per the task's order of work, Section 1 is
% written last, once Sections 2-9 are settled. Do not draft contribution
% claims here until the Results/Discussion sections determine what each
% claim can honestly say.

% ------------------------------------------------------------------
\section{Related Work}
\label{sec:related-work}

The work presented in this paper intersects six established research
paradigms: automated machine learning \citep{hutter2019automl,elsken2019nas,feurer2022autosklearn2,he2021automlsurvey}, information
extraction \citep{banko2007openie,lu2022uie,josifoski2022genie,liu2024openiesurvey}, retrieval-augmented generation \citep{guu2020realm,lewis2020rag}, parameter-efficient
fine-tuning \citep{houlsby2019peft,lester2021prompttuning,hu2021lora,dettmers2023qlora}, feedback-based self-improvement \citep{wang2023selfinstruct,madaan2023selfrefine,gulcehre2023rest,chen2024spin}, and probabilistic knowledge
representation \citep{bach2017psl,ji2021kgsurvey,nguyen2020kgcsurvey}. Each paradigm addresses a partial subset of the
capabilities that SLDE-AFT integrates within a single architecture; the
subsections below review each in turn, and Section~\ref{sec:related-work-dysect}
discusses the one system in the literature that spans the largest share of
these capabilities at once.

\subsection{Automated Machine Learning}

Automated machine learning (AutoML) emerged to reduce human effort in
constructing supervised learning pipelines by automating model selection,
hyperparameter optimisation, and architecture search \citep{hutter2019automl,he2021automlsurvey}.
Bayesian-optimisation-based frameworks such as Auto-Sklearn \citep{feurer2015autosklearn}
and its successor Auto-Sklearn~2.0 \citep{feurer2022autosklearn2} introduced
meta-learning-guided warm starting and portfolio construction to reduce
pipeline search time substantially. Ensemble-driven systems such as
AutoGluon \citep{erickson2020autogluon} extended automation to tabular
prediction tasks through multi-layer stacking, while commercial platforms
such as Amazon SageMaker Autopilot \citep{das2020sagemaker} enabled
transparent, editable pipeline generation at deployment scale. Neural
architecture search methods \citep{elsken2019nas} further widened the
automation boundary to model structure itself. Despite these advances, AutoML
frameworks treat data extraction as a fixed upstream preprocessing step: the
adaptive boundary begins at feature engineering and ends at model deployment,
leaving the connection between heterogeneous source ingestion, knowledge
accumulation, and iterative model adaptation outside the automation scope
addressed by the closed-loop design described in Section~4.

\subsection{Information Extraction and Relation Extraction}

Information extraction research has pursued two broad trajectories:
schema-free open extraction and schema-guided structured extraction. Open
Information Extraction systems \citep{banko2007openie,liu2024openiesurvey}
showed that relational triples could be extracted from large text corpora
without predefining relation types, using lightweight linguistic processing
and redundancy-based confidence estimation. Document-level relation
extraction benchmarks such as DocRED \citep{yao2019docred} established
evaluation standards for reasoning across sentence boundaries, while
supervised relation classification datasets such as TACRED
\citep{zhang2017tacred} became widely adopted benchmarks for sentence-level
relation extraction. Generative extraction architectures reformulated triple
extraction as sequence-to-structure generation: REBEL \citep{cabot2021rebel}
reformulated end-to-end relation extraction as autoregressive triplet
generation, GenIE \citep{josifoski2022genie} extended this to
schema-constrained decoding at large entity-relation scale, and TANL
\citep{paolini2021tanl} treated structured prediction as translation between
augmented natural languages. Universal IE frameworks \citep{lu2022uie,wang2023instructuie}
unified heterogeneous extraction tasks under a single architecture via
structured extraction languages and multi-task instruction tuning. Joint
extraction systems such as DyGIE++ \citep{wadden2019dygiepp} improved the
integration of entity and relation extraction, while open-world attribute
discovery systems such as OA-Mine \citep{zhang2022oamine} enabled discovery
of previously unseen attributes beyond predefined schemas. Zero-shot relation
extraction approaches \citep{levy2017zeroshot} demonstrated that
question-template reformulation can facilitate extraction of previously
unseen relation types without retraining. These approaches improve
extraction accuracy, coverage, schema flexibility, or generalisation
capability, but they generally treat extraction as a terminal process and do
not leverage accumulated extracted knowledge to drive iterative model
adaptation through an automated feedback loop.

\subsection{Retrieval-Augmented Generation}

Retrieval-augmented generation (RAG) addresses the knowledge limitations of
parametric language models by augmenting generation with dynamically
retrieved external evidence, reducing hallucination and improving factual
grounding. The RAG framework of \citet{lewis2020rag} demonstrated that dense
passage retrieval combined with sequence-to-sequence generation improves
factual accuracy on knowledge-intensive question-answering and reasoning
tasks. REALM \citep{guu2020realm} integrated retrieval directly into language
model pre-training, enabling differentiable retrieval learning while
decoupling parametric knowledge from external evidence sources. SELF-RAG
\citep{asai2023selfrag} introduced adaptive retrieval and self-reflection
tokens that let the model critique its own outputs and selectively retrieve
evidence, improving factual attribution. HippoRAG \citep{gutierrez2024hipporag}
proposed a neurobiologically inspired long-term memory architecture combining
knowledge graph construction with personalised PageRank retrieval, improving
multi-hop reasoning. GraphRAG \citep{edge2025graphrag} extended RAG toward
graph-structured community summarisation, enabling retrieval over
hierarchically organised entity relationships. Despite these advances,
existing RAG frameworks primarily focus on retrieval-enhanced inference and
generally treat retrieved knowledge as a transient resource used only during
generation; accumulated knowledge is not typically used to drive iterative
extraction improvement, automated supervision generation, or continuous
model adaptation.

\subsection{Parameter-Efficient Fine-Tuning}

The pretrain-then-fine-tune paradigm established by BERT \citep{devlin2019bert}
demonstrated that large-scale bidirectional pretraining followed by
task-specific adaptation can achieve strong downstream NLP performance;
however, adaptation typically relies on task-specific labelled datasets.
Instruction tuning via FLAN \citep{wei2022flan} and RLHF-aligned adaptation
via InstructGPT \citep{ouyang2022instructgpt} showed that natural-language
task descriptions and preference signals can generalise fine-tuned behaviour
to unseen tasks, yet both approaches depend on large-scale human annotation.
Parameter-efficient fine-tuning (PEFT) methods substantially reduced
adaptation costs by updating only a small subset of model parameters.
Representative approaches include adapter layers \citep{houlsby2019peft},
soft prompt tuning \citep{lester2021prompttuning}, LoRA \citep{hu2021lora},
and QLoRA \citep{dettmers2023qlora}, which progressively improved parameter
efficiency while preserving downstream performance. Direct Preference
Optimization \citep{rafailov2023dpo} replaced reward modelling with a
simpler binary classification objective. Although these approaches
significantly reduce adaptation overhead, they generally rely on externally
curated supervision and are initiated through manually defined training
workflows; the automatic generation of supervision signals from accumulated
extracted knowledge, and the continuous triggering of model adaptation from
that signal, remain largely unexplored by this line of work.

\subsection{Feedback-Based and Self-Improving Learning}

A growing line of research investigates how language models can improve
their own outputs through iterative self-generated feedback, without
additional human supervision. SELF-REFINE \citep{madaan2023selfrefine}
showed that a single model acting as both generator and critic can improve
output quality through training-free iterative refinement. Reinforced
Self-Training \citep{gulcehre2023rest} proposed a Grow-Improve loop in which
the model generates candidate outputs, filters them by reward score, and
retrains on high-quality examples offline, reducing the cost of RLHF.
Self-Play Fine-Tuning \citep{chen2024spin} enabled iterative self-improvement
by training a model to distinguish its own outputs from human-quality
references, without additional human data. SELF-INSTRUCT
\citep{wang2023selfinstruct} introduced a self-bootstrapping pipeline for
synthetic instruction generation that reduces annotation dependency.
Toolformer \citep{schick2023toolformer} showed that language models can
acquire external tool-use capabilities through self-supervised data
generation. Continual learning research \citep{wang2024continuallearning}
and Learning without Forgetting \citep{li2017lwf} address stability-plasticity
trade-offs in sequential adaptation settings. These approaches improve
generation quality, self-supervision, or adaptation efficiency, but they
generally operate on transient model outputs and do not maintain a
persistent structured knowledge repository that accumulates extracted facts
and converts them into supervision signals for continuous model adaptation.

\subsection{Probabilistic Knowledge Representation and Explainability}

Principled uncertainty modelling over extracted facts matters for systems
that must decide which knowledge to trust, accumulate, and reuse for
downstream adaptation. Probabilistic Soft Logic \citep{bach2017psl} provided
a scalable framework for soft reasoning over structured relational domains
using hinge-loss Markov random fields, showing that confidence-weighted
inference over structured variables is both theoretically grounded and
practically efficient. Knowledge graph surveys \citep{ji2021kgsurvey,nguyen2020kgcsurvey}
identify confidence-aware extraction, knowledge base population, and
uncertainty management as persistent challenges in scalable knowledge-centric
systems. Explainability research has established principled attribution
methods: LIME \citep{ribeiro2016lime} provides locally faithful
model-agnostic explanations through input perturbation, and LLM
explainability surveys \citep{zhao2023explainability} extend attribution
taxonomies to the prompting and fine-tuning paradigms of large language
models. These works show that confidence modelling and attribution are
tractable and important, but they operate primarily as post-hoc or offline
analysis tools; the integration of provenance traceability and
confidence-aware filtering as \emph{active} components of a continuous
extraction and adaptation pipeline is comparatively underexplored.

\subsection{Closest Prior System: DySECT}
\label{sec:related-work-dysect}

The system closest to SLDE-AFT in the published literature is DySECT, a
Dynamic Self-Evolving Extraction and Curation Toolkit, presented in a paper
titled ``A Dynamic Self-Evolving Extraction System''
\citep{aminnaseri2026dysect}. DySECT incrementally populates a self-expanding
knowledge base with triples extracted by an LLM, aggregates them with a
probabilistic confidence and graph-based reasoning layer built on top of
Theo \citep{mitchell1989theo}, and feeds the enriched knowledge base back
into the extractor through prompt tuning, few-shot example sampling, or
fine-tuning on knowledge-base-derived synthetic data -- forming, in the
authors' own description, a closed loop in which extraction improves
knowledge and knowledge improves extraction. Table~1 of
\citet{aminnaseri2026dysect} reports that, at the first KB-guided iteration,
recall improves by 5--8 percentage points over the no-feedback baseline
across all four evaluated models (GPT-4.1: 22.80\%$\to$30.62\%; GPT-4.1-mini:
12.72\%$\to$20.22\%; Kimi~K2.5: 32.03\%$\to$37.85\%; LLaMA-3.3~70B:
23.23\%$\to$28.78\%, positive-feedback mode), on DocRED. DySECT also keeps
its knowledge base explicitly inspectable, letting a user review triples,
confidence statistics, and hierarchical structure, and manually validate,
invalidate, or add knowledge.

\paragraph{The confidence aggregation formula is not merely similar --
it is the same formula.} \citet{aminnaseri2026dysect}'s Equations 1--2 define
$C_{\mathrm{agg}}(t) = 1 - \prod_i (1-\lambda c_i)^{f_i}$ with a default
shrinkage $\lambda=0.75$, followed by a mutual-exclusivity penalty
$C(t) = C_{\mathrm{agg}}(t) / (m(t)+1)$. This is notationally and
mathematically identical to SLDE-AFT's Conservative Noisy-OR aggregation
with mutual-exclusivity penalty (Section~5; the same default shrinkage of
0.75 is independently confirmed in this project's own
\texttt{scripts/dec003\_product\_probkb\_run.py}). This paper does not claim
priority or originality for this formula, and does not claim it is novel or
absent from prior extraction-KB systems: DySECT is a published, citable
counter-example to any such claim. What Section~5 contributes instead is
a formal derivation and proofs of this formula's boundedness, monotonicity,
and saturation properties, and the rival-ceiling result
$C(t) < 1/(m(t)+1) \le 0.5$ for any contested single-valued-predicate slot
-- properties stated but not proved in \citet{aminnaseri2026dysect}'s
description of the same formula.

\paragraph{Capabilities SLDE-AFT shares with DySECT.} Four elements of
DySECT's design, as described by \citet{aminnaseri2026dysect}, are also
present and verified in SLDE-AFT's implementation (Section~4): a
\emph{closed loop} in which extraction populates a knowledge base and the
knowledge base feeds back into extraction; a \emph{probabilistic confidence}
layer over accumulated triples, rather than treating every extraction as
equally trustworthy; \emph{prompt-augmentation feedback}, in which the
current state of the knowledge base is used to steer the extractor's next
pass (DySECT via prompt tuning and few-shot example sampling; SLDE-AFT via
the Feedback Controller's Encouraging/Prohibitive-mode prompt updates,
Section~4); and \emph{synthetic training data generated from high-confidence
knowledge-base triples} to fine-tune the extractor, rather than from an
externally curated dataset.

\paragraph{Capabilities DySECT reports that SLDE-AFT's implementation does
not have.} \citet{aminnaseri2026dysect} describe two further capabilities
that were also specified in SLDE-AFT's original architecture design but are
not present in the implementation evaluated in this paper, verified directly
against the current codebase: \emph{hierarchical knowledge abstraction}
(DySECT accumulates and exposes hierarchical domain concept structure;
SLDE-AFT's design document specified KNN-clustering-based subconcept
abstraction, but no such clustering or hierarchy-construction code exists in
the evaluated system) and a \emph{human curation interface} allowing a user
to inspect, validate, invalidate, or add knowledge-base entries directly
(SLDE-AFT's design document specified an analogous human validation
interface with character-level source offsets and a model-version registry
with rollback; none of these are implemented in the evaluated system).
These are stated here as an honest scope limitation of the current
implementation relative to DySECT, not as a design decision, and are
revisited in Section~9 (Limitations).

\paragraph{Capabilities unique to SLDE-AFT, verified as implemented.} Two
elements of SLDE-AFT's design have no counterpart described in
\citet{aminnaseri2026dysect}. First, dual-source structured seeding: SLDE-AFT
ingests structured records (CSV/JSON) through a deterministic schema-based
extraction path at a fixed high confidence, in addition to the LLM-based
unstructured path, before both sources are aggregated by the same
probabilistic mechanism (Section~4); DySECT's knowledge base is populated
from LLM extraction alone. Second, provenance as an automated training-data
gate: rather than a human curation step, SLDE-AFT's provenance layer
automatically withholds a triple from becoming fine-tuning training data
unless it has structured corroboration (Section~4, Section~7) -- an
automated mechanism aimed at synthetic-data quality, verified to measurably
change training-data precision on real data (Section~7), not a substitute
for DySECT's manual interface but a different kind of gate entirely. The
confidence aggregation formula itself is not unique to SLDE-AFT, as stated
above; the formal analysis of that formula (Section~5) is what this paper
contributes on top of it.

\paragraph{Empirical comparability.} The two systems' empirical evaluations
use different benchmarks, base models, and metrics: DySECT reports a
5--8 percentage point recall lift from KB-guided extraction with no
retraining, and recall improvements on DocRED across GPT-4.1, GPT-4.1-mini,
LLaMA-3.3~70B, and Kimi~K2.5; SLDE-AFT's evaluation (Sections~6--8) uses
CaRB, DocRED, and BioRED with a different set of base models and a
seed-based statistical protocol. No direct numerical comparison between the
two systems is made in this paper as a result; a controlled head-to-head
evaluation on a shared benchmark is identified as future work in
Section~10.

\subsection{Synthesis}

Across the paradigms reviewed above, a consistent pattern emerges: each
research tradition addresses one layer of the extraction-knowledge-adaptation
problem in relative isolation. AutoML research automates model selection and
optimisation while leaving knowledge acquisition outside the automation
boundary \citep{hutter2019automl,feurer2022autosklearn2}. Information
extraction systems improve fact discovery but generally treat extracted
knowledge as a terminal output rather than a reusable learning signal
\citep{banko2007openie,lu2022uie,liu2024openiesurvey}. RAG improves grounding
but keeps the retriever static \citep{lewis2020rag,asai2023selfrag}.
Fine-tuning reduces adaptation cost but requires manually curated
supervision and manually triggered training \citep{hu2021lora,dettmers2023qlora}.
Self-improving systems refine outputs but lack a structured, persistent
knowledge substrate \citep{madaan2023selfrefine,gulcehre2023rest}. Provenance
and confidence-modelling systems record attribution but do not typically use
it as an active training-data filter \citep{bach2017psl,ribeiro2016lime}.
DySECT \citep{aminnaseri2026dysect} is the one system in this review that
combines several of these layers into a closed loop, and the preceding
subsection states plainly where SLDE-AFT's design overlaps with it and where
it differs. Section~3 formalises the resulting cross-paradigm gap and
positions SLDE-AFT's architecture against it.

% ------------------------------------------------------------------
\section{Cross-Paradigm Research Gap}
% TODO: not yet written.

\section{Proposed Framework}
% TODO: not yet written.

\section{Formal Analysis of the Aggregation Rule}
\label{sec:formal-analysis}

\subsection{Definition and Attribution}
\label{sec:formal-analysis-definition}

As established in Section~\ref{sec:related-work-dysect}, the confidence
aggregation rule used throughout this paper is not original to it:
\citet{aminnaseri2026dysect}'s Equations 1--2 define the identical
formula, under the identical default shrinkage. This paper cites that
formula rather than claiming it. For a functionally single-valued
predicate slot with candidate object $t$ supported by observations with
self-reported confidences $c_1,\dots,c_n \in [0,1]$
% [EVID-041, "confidence is self-reported, not derived from calibration
% or token probabilities" -- verified against src/extractors/openrouter_llm.py]
and shrinkage $\lambda \in (0,1]$ (default $\lambda=0.75$
% [scripts/dec003_product_probkb_run.py, scripts/dec006_scaleup_probkb_run.py --
% independently confirmed default, matching aminnaseri2026dysect Eq. 1]
), the \emph{conservative Noisy-OR support} is
\begin{equation}
A(t) \;=\; 1 - \prod_{i=1}^{n} \bigl(1 - \lambda c_i\bigr).
\label{eq:noisyor}
\end{equation}
For a slot designated functionally single-valued, with $m(t)$ competing
object values already accepted for the same (subject, predicate) slot,
the admitted confidence is
\begin{equation}
C(t) \;=\; \frac{A(t)}{m(t) + 1}.
\label{eq:conflict}
\end{equation}
For a non-functional predicate (only \texttt{has\_color} in the
product-domain configuration
% [configs/functional_predicates_product_domain.json]
), $C(t) = A(t)$ unconditionally: no mutual-exclusivity penalty is
applied, since multiple simultaneous values are legitimate. A candidate
is admitted into the knowledge base when $C(t) \ge \tau$, with
$\tau=0.88$ the project's standing operating threshold
% [scripts/dec003_product_probkb_run.py, scripts/dec006_scaleup_probkb_run.py]
throughout the experiments in Sections~\ref{sec:experimental-setup}--\ref{sec:results}.
What this paper contributes on top of Equations~\ref{eq:noisyor}--\ref{eq:conflict}
is (i) formal proofs of properties stated but not proved in
\citet{aminnaseri2026dysect}'s presentation of the same rule, and (ii) an
empirical test, on two independent real-data snapshots, of whether a
corrected rule can outperform it. Both are given below.

\subsection{Formal Properties}
\label{sec:formal-analysis-properties}

\begin{proposition}[Range and boundedness]
For any non-empty finite multiset of confidences $c_1,\dots,c_n \in [0,1]$
and $\lambda \in (0,1]$, $A(t) \in [0,1)$ whenever at least one $c_i < 1$,
and $A(t) \to 1$ only in the limit $n\to\infty$ with every $c_i$ bounded
away from $0$. Consequently $C(t) \in [0,1)$ for every functionally
single-valued slot with $m(t) \ge 0$.
\end{proposition}
\begin{proof}
Each factor $(1-\lambda c_i) \in (0,1]$ since $\lambda c_i \in [0,1]$, so
the product $\prod_i(1-\lambda c_i) \in (0,1]$ and is strictly positive
for finite $n$; hence $A(t) = 1 - \prod_i(1-\lambda c_i) < 1$ strictly for
every finite observation set. $A(t) \ge 0$ follows since each factor is
$\le 1$, so the product is $\le 1$. Dividing by $m(t)+1 \ge 1$ preserves
$C(t) \in [0,1)$.
\end{proof}

\begin{proposition}[Order-invariance]
$A(t)$ depends only on the multiset $\{c_1,\dots,c_n\}$, not on the order
in which observations are accumulated.
\end{proposition}
\begin{proof}
Multiplication is commutative, so $\prod_i(1-\lambda c_i)$ is invariant
under any permutation of the index set $i=1,\dots,n$.
\end{proof}

\begin{proposition}[Monotonicity in evidence]
\label{prop:monotone}
Adding an observation with confidence $c_{n+1} > 0$ to an existing set of
$n$ observations strictly increases $A(t)$; $A(t)$ is non-decreasing in
$n$ and strictly increasing whenever every added observation has
$c_i > 0$.
\end{proposition}
\begin{proof}
Let $A_n = 1-\prod_{i=1}^n(1-\lambda c_i)$. Then
$A_{n+1} - A_n = \prod_{i=1}^n(1-\lambda c_i) \cdot \lambda c_{n+1} > 0$
whenever $c_{n+1} > 0$, since $\prod_{i=1}^n(1-\lambda c_i) > 0$ by the
boundedness proposition above.
\end{proof}

\begin{proposition}[Saturation under repeated observation]
\label{prop:saturation}
If the same confidence $c$ is observed $n$ times, $A_n = 1-(1-\lambda c)^n
\to 1$ as $n \to \infty$, geometrically, at rate $(1-\lambda c)$: the
residual $1 - A_n = (1-\lambda c)^n$ decreases by the constant factor
$(1-\lambda c)$ with every additional identical observation.
\end{proposition}
\begin{proof}
Immediate from Equation~\ref{eq:noisyor} with $c_i=c$ for all $i$.
\end{proof}

\begin{proposition}[Corroboration requirement]
\label{prop:corroboration}
With $\lambda=0.75$ and $\tau=0.88$, no single observation, regardless of
its confidence, is sufficient for admission: a triple requires at least
two corroborating observations to reach $\tau$.
\end{proposition}
\begin{proof}
A single observation with the maximum possible confidence $c_1=1$ yields
$A(t) = \lambda = 0.75 < \tau = 0.88$ from Equation~\ref{eq:noisyor} with
$n=1$. More generally, admission with $n$ observations each at $c=1$
requires $1-(1-\lambda)^n \ge \tau$, i.e.\ $n \ge \ln(1-\tau)/\ln(1-\lambda)
= \ln(0.12)/\ln(0.25) \approx 1.53$, so $n \ge 2$.
\end{proof}
This property is visible directly in the project's own data: a
functionally single-valued triple supported by exactly one observation
(the common case for a structured-only fact, fixed confidence $0.98$
% [EVID-041, src/extractors/structured.py -- structured confidence
% fixed at 0.98]
) scores $A(t) = \lambda \times 0.98 = 0.735$ under
Equation~\ref{eq:noisyor} -- below $\tau=0.88$ regardless of how
trustworthy that single structured source is, exactly as
Proposition~\ref{prop:corroboration} predicts.

\begin{proposition}[Conflict-penalty properties]
\label{prop:conflict}
For fixed $A(t)$, $C(t) = A(t)/(m(t)+1)$ is strictly decreasing in
$m(t)$, and $C(t) = A(t)$ exactly when $m(t)=0$ (no penalty when the slot
is uncontested).
\end{proposition}
\begin{proof}
$\partial C(t)/\partial m(t) = -A(t)/(m(t)+1)^2 < 0$ for $A(t) > 0$, and
$C(t)|_{m(t)=0} = A(t)/1 = A(t)$ directly from Equation~\ref{eq:conflict}.
\end{proof}

\begin{proposition}[Rival ceiling]
\label{prop:ceiling}
For any functionally single-valued predicate slot with $m(t) \ge 1$
competing object values, $C(t) < 1/(m(t)+1) \le 0.5$, independent of how
much observational support $A(t)$ accumulates.
\end{proposition}
\begin{proof}
By the boundedness proposition, $A(t) < 1$ for any finite observation
set. Hence $C(t) = A(t)/(m(t)+1) < 1/(m(t)+1)$. For $m(t) \ge 1$,
$1/(m(t)+1) \le 1/2$.
\end{proof}

\begin{corollary}
\label{cor:ceiling-consequence}
With $\tau = 0.88 > 0.5$, a contested functionally single-valued slot
($m(t) \ge 1$) can never be admitted under Equations~\ref{eq:noisyor}--\ref{eq:conflict},
regardless of how many corroborating observations accumulate for the
leading candidate. Combined with the fact that no candidate is ever
deleted or pruned from the accepted-candidates store anywhere in
\texttt{src/probkb\_v2\_adapter.py} or \texttt{src/pkb\_instrumentation.py}
% [Results Summary.md, Claim 3 -- verified directly, no del/pop/remove
% on the accepted dict]
, a genuinely contested functional slot is not merely slow to resolve
under this rule: it is structurally unresolvable at $\tau=0.88$.
\end{corollary}

\begin{proposition}[Independence assumption and repeated observation]
\label{prop:independence}
Equation~\ref{eq:noisyor} is a Noisy-OR combination rule, which assumes
each $c_i$ is drawn from a conditionally independent evidence source.
This assumption is violated when the same extractor, invoked at
temperature $0$, is asked to re-read the same source text: the output is
deterministic, so repeated invocations produce byte-identical
observations that are treated by Equation~\ref{eq:noisyor} as though they
were independent corroboration, inflating $A(t)$ by
Proposition~\ref{prop:saturation}'s saturation dynamic regardless of
whether the underlying claim is correct.
\end{proposition}
This is not merely a theoretical caveat: 31 unstructured-only triples in
the DEC-018 200-product snapshot were observed 2--21 times each by
exactly this mechanism, all 31 incorrect against gold
% [EVID-029]
, and the same snapshot's real-data expected calibration error is
$0.3332$
% [EVID-036]
, both consistent with repeated non-independent observation being
mistaken for corroboration. Section~\ref{sec:formal-analysis-empirical}
tests a corrected rule that aggregates over distinct sources instead of
repeat observations, directly targeting this violated assumption.

\subsection{Computational Complexity}
\label{sec:formal-analysis-complexity}

The production candidate-acceptance path is empirically $O(N)$ per call
and $O(N^2)$ cumulative over $N$ accumulated observations
% [EVID-036]
, caused by a linear scan over every previously accepted candidate in
\texttt{src/pkb\_instrumentation.accepted\_slot\_keys} on every new
observation -- measured at $72.7\times$ latency growth for $80\times$
more observations
% [EVID-036]
. An indexed alternative (same aggregation mathematics, keyed by
normalized (subject, predicate) slot rather than scanned linearly)
achieves $O(1)$ per call and $O(N)$ cumulative
% [EVID-036]
; it was implemented illustratively but not wired into the evaluated
system. This is not a limiting factor at the scale evaluated in this
paper ($N$ in the hundreds, Section~\ref{sec:experimental-setup}), where
LLM API latency dominates wall-clock time
% [EVID-036, cross-referenced against DEC-008's linear end-to-end runtime
% finding]
, but is stated here as a forward-looking scalability property of the
current implementation, separate from the aggregation rule's own
mathematical properties above.

\subsection{Empirical Validation: The Ceiling, Tested Corrections, and Their Results}
\label{sec:formal-analysis-empirical}

Corollary~\ref{cor:ceiling-consequence} is a proof, not a measurement: it
establishes that the ceiling \emph{must} block every contested
single-valued slot, but says nothing about whether contested slots
actually arise on real data, nor whether a corrected rule can do better.
Both questions were tested directly
% [DEC-026, EVID-041]
, offline, by replaying stored observations from two independent
gold-labeled snapshots -- 657 triples from 35 products
% [EVID-013, EVID-014, EVID-036]
and 2241 triples from 140 products
% [EVID-029, EVID-030]
-- through six scoring rules applied to the identical observations (only
the aggregation formula varies), each with its own validation-selected
threshold and a subject-level (product-level) held-out test split (35
and 92 test-split subjects respectively), plus every rule re-evaluated
at the fixed pipeline threshold $\tau=0.88$
% [EVID-041]
.

\paragraph{The ceiling is confirmed empirically, at both thresholds
tested.} The published rule (Equations~\ref{eq:noisyor}--\ref{eq:conflict})
admits exactly zero contested single-valued-predicate slots on either
test split, at either its own F1-selected threshold or the fixed
$\tau=0.88$
% [EVID-041]
. An unconstrained max-merge baseline ($C(t)=\max_i c_i$, no conflict
handling), evaluated identically, admits 6 and 40 contested slots at its
own F1-selected threshold and 20 and 69 at $\tau=0.88$, on the two
datasets respectively
% [EVID-041]
-- confirming that it is the ceiling, and not an absence of contested
slots in the data, that accounts for the published rule's zero.

\paragraph{Tested corrections that are not claimed as contributions.} Two
candidate corrections to Equation~\ref{eq:conflict} were tested and are
reported here as null or negative results, not as improvements:
\begin{itemize}
  \item An evidence-share reformulation, $C(t) = A(t)^2 / \sum_j A(t_j)$
  summed over every candidate $j$ in the same slot, removes the
  algebraic ceiling of Proposition~\ref{prop:ceiling} (it reduces
  correctly to $C(t)=A(t)$ when $m(t)=0$, matching
  Proposition~\ref{prop:conflict}) but is an exact empirical null: it
  produces an admitted set identical to Equation~\ref{eq:conflict} on
  both datasets (bootstrap 95\% CI for the $F_1$ difference exactly
  $[0.0000,0.0000]$ over $10{,}000$ subject resamples)
  % [EVID-041]
  , because no observed contested-slot score in either dataset approaches
  the operating threshold (maximum observed score $0.581$ against
  $\tau=0.73$ on the 657-triple dataset)
  % [EVID-041]
  .
  \item An unsquared variant, $C(t) = A(t)/\sum_j A(t_j)$, was tested
  afterward as an exploratory, post-hoc check and is a negative result:
  it degenerates to $C(t)=1$ for every uncontested slot ($m(t)=0
  \Rightarrow \sum_j A(t_j) = A(t) \Rightarrow C(t)=1$ regardless of the
  actual evidence strength), discarding real signal on $495$ of $495$
  uncontested functional-predicate rows in the 657-triple dataset
  % [EVID-041]
  , and its $F_1$ is significantly worse than the published rule on one
  dataset (95\% CI $[-0.0870,-0.0118]$) and not significantly different
  on the other
  % [EVID-041]
  . This failure is itself a formal, not merely empirical, result: any
  share-based replacement for Equation~\ref{eq:conflict} must satisfy
  $C(t)=A(t)$ when $m(t)=0$ to avoid this degeneracy -- the squared form
  above satisfies this reduction and the unsquared form does not, which
  is why one is empirically inert and the other is empirically harmful,
  independent of which snapshot either is evaluated on
  % [EVID-041]
  .
\end{itemize}

\paragraph{The corrected rule: distinct-source aggregation.} The one
tested correction that is claimed as a contribution targets
Proposition~\ref{prop:independence}'s violated-independence source of
error directly, rather than the ceiling: aggregate
Equation~\ref{eq:noisyor} over confidences from distinct
$(\text{source\_type}, \text{source\_id})$ pairs (deduplicating repeated
observations from the same source, keeping the maximum confidence per
source) instead of over every repeat observation, leaving
Equation~\ref{eq:conflict}'s conflict handling otherwise unchanged. At
the pipeline's fixed operating threshold $\tau=0.88$ -- the threshold
this paper's other results use, and therefore the relevant comparison --
this correction yields a statistically significant precision gain over
the published rule at \emph{exactly unchanged} recall on both datasets
(bootstrap 95\% CI for the recall difference exactly $[0.0000,0.0000]$
on both):
\begin{center}
\begin{tabular}{@{}lccc@{}}
\toprule
Dataset & Precision diff.\ (95\% CI) & Recall diff.\ (95\% CI) & $F_1$ diff.\ (95\% CI) \\
\midrule
657 triples / 35 products & $+0.1060\ [0.0321, 0.2267]$ & $0.0000\ [0.0000, 0.0000]$ & $+0.0192\ [0.0057, 0.0395]$ \\
2241 triples / 140 products & $+0.0270\ [0.0076, 0.0519]$ & $0.0000\ [0.0000, 0.0000]$ & $+0.0059\ [0.0016, 0.0115]$ \\
\bottomrule
\end{tabular}
\end{center}
% [EVID-041, tau=0.88 fixed-threshold bootstrap, 10,000 resamples over test-split subjects]
Both the precision and $F_1$ intervals exclude zero on both independent
datasets. The same correction also improves $F_1$ at each dataset's own
F1-selected threshold (+0.0106 $[0.0028,0.0222]$ and +0.0039
$[0.0013,0.0075]$ respectively)
% [EVID-041]
, so the effect is not an artifact of threshold selection. An
unconstrained max-merge baseline also outperforms the published rule at
its own F1-selected threshold, matching the same direction already found
in this project's $N{=}50$ closed-loop ablation
% [EVID-039, AUDIT.md A1]
, but that advantage does not survive at the shared $\tau=0.88$ operating
threshold (bootstrap $F_1$-difference CI includes zero on both datasets,
driven by a significant precision loss that offsets a significant recall
gain)
% [EVID-041]
and is not reported as a corrected rule here, only as further empirical
confirmation of Proposition~\ref{prop:ceiling} (Section~\ref{sec:formal-analysis-empirical},
above).

\subsection{Summary}
\label{sec:formal-analysis-summary}

Three claims follow from this section. First, the aggregation formula
(Equations~\ref{eq:noisyor}--\ref{eq:conflict}) and its default shrinkage
$\lambda=0.75$ are \citet{aminnaseri2026dysect}'s, cited here, not
claimed as this paper's contribution (Section~\ref{sec:related-work-dysect}).
Second, this paper's contribution is the formal analysis above --
boundedness, order-invariance, monotonicity, the corroboration
requirement, saturation, the conflict-penalty properties, and the rival
ceiling with its structural-unresolvability corollary -- together with
the empirical demonstration that this ceiling blocks every contested
single-valued slot on two independent real-data snapshots, at both the
F1-selected and the fixed pipeline threshold. Third, the corrected rule
this paper proposes is distinct-source aggregation
(Section~\ref{sec:formal-analysis-empirical}), reported with its
measured precision and $F_1$ gain over the published rule at
$\tau=0.88$ and the corresponding 95\% confidence intervals above; the
two other candidate corrections tested (evidence-share and its unsquared
variant) are reported as tested and null or negative, not as
contributions of this paper.

\section{Experimental Setup}
\label{sec:experimental-setup}
% TODO: not yet written.

\section{Results}
\label{sec:results}
% TODO: not yet written -- awaiting the user's Section 7 Word document
% (6 tables, 4 figures) and its supplementary material (Tables S1-S8),
% per the task's "Existing material to reconcile against" instruction.

\section{Discussion}
% TODO: not yet written. Carry forward from Section 2.7 / Learnings.md:
% DySECT reports consistent recall IMPROVEMENTS from KB-guided extraction
% on DocRED (Table 1, all four models), whereas SLDE-AFT's own DocRED
% pilot (DEC-020) found the PKB aggregation mechanism REDUCED recall
% (naive baseline 0.032 -> PKB-aggregated 0.005). State this contrast
% honestly here -- do not write it yet, per the task's staged order.

\section{Limitations}
% TODO: not yet written.

\section{Conclusion}
% TODO: not yet written.

\section*{Declarations}
% TODO: not yet written -- data/code availability, funding, competing
% interests, author contributions. Leave TODO markers for the user to
% supply text where required.

\bibliographystyle{unsrtnat}
\bibliography{references}

\end{document}
